\section{Conventions and elementary tools}\label{sec:pre}

This section fixes the conventions of both documents and proves the elementary lemmas that the
other sections use without comment: the flat step and the flat concatenation of paths and isotopies,
cutoffs, the top power and the local form of toric symplectic forms, and the Moser lemma with the sign
conventions of the paper. Lemma~\ref{lem:pre-top-power} contains \Mref{Lemma}{lem:hdl-top-power} as a
special case. Tables~\ref{tab:pre-conventions} and~\ref{tab:pre-interfaces} (\S\ref{sec:pre-unified}) collect the
coordinate, sign and normalisation conventions of all parts of the construction and say where one is converted into
another; Table~\ref{tab:pre-notation} extends \Mref{Table}{tab:mech-notation}.

\subsection{Conventions}\label{sec:pre-conventions}

This subsection expands the conventions of \Msec{sec:intro-conventions} with the identities used in
the proofs.

Manifolds are finite dimensional, Hausdorff and second countable; boundaries are stated explicitly.

\emph{Phase space.} We write $\R^{2n}=\C^n$ with $z_j=x_j+iy_j$ and
$\omega_0=\sum_{j=1}^ndx_j\wedge dy_j$. Let $J_0$ be the matrix of $\omega_0$, i.e.\
$\omega_0(v,w)=\langle v,J_0w\rangle$; in the coordinates $(x,y)$, $J_0=\bigl(\begin{smallmatrix}0&I\\-I&0
\end{smallmatrix}\bigr)$, which is multiplication by $-i$ on $\C^n$. Where $\R^{2n}$ is regarded as a
complex vector space with $i$ acting by $J_0$ (\Msec{sec:ham}), this complex structure is the conjugate of
that of $\C^n$. For a function $H(t,z)$, smooth in $(t,z)$, the Hamiltonian vector field $X_H=X_{H_t}$ is
defined by $\iota_{X_H}\omega_0=dH_t$, where $d$ is the spatial differential. Equivalently
\[
 X_H=J_0\nabla H,\qquad \dot x_j=\frac{\partial H}{\partial y_j},\quad \dot y_j=-\frac{\partial H}{\partial x_j},
 \qquad \dot z_j=-2i\,\frac{\partial H}{\partial\bar z_j},
\]
since $\omega_0(X,w)=\langle X,J_0w\rangle=-\langle J_0X,w\rangle$ and $J_0^2=-I$. We write $\varphi^t_H$ for the
flow of $X_H$ from time $0$ to time $t$. For example $H=\pi|z_1|^2$ generates
$z_1\mapsto e^{-2\pi it}z_1$, a rotation of period one. The letter $p$ is reserved for moment maps.

\emph{Capacities and balls.} For $c>0$ and $s\in\C^n$,
\[
 B(c)=\{z\in\C^n:\pi|z|^2\le c\},\qquad \Bo(c)=\{\pi|z|^2<c\},\qquad B(s,c)=s+B(c),
\]
the closed and open balls of capacity $c$; the Euclidean radius of $B(c)$ is $\sqrt{c/\pi}$. We write
$\EB(r)$ and $\EB^\circ(r)$ for the closed and open Euclidean balls of radius $r$ about the origin.
Volumes are taken with respect to $\omega^n/n!$; for $\omega_0$ this is Lebesgue measure, and
$\vol B(c)=\pi^n(c/\pi)^n/n!=c^n/n!$.

\begin{lemma}[Dilations]\label{lem:pre-dilation}
Let $\lambda>0$, $\delta_\lambda(z)=\lambda z$, and let $H(t,z)$ be smooth. Then
$\delta_\lambda^*\omega_0=\lambda^2\omega_0$, $\delta_\lambda(B(s,c))=B(\lambda s,\lambda^2c)$, and the function
$\tilde H(t,z)=\lambda^2H(t,z/\lambda)$ satisfies $X_{\tilde H}(t,z)=\lambda X_H(t,z/\lambda)$; hence an integral curve
$\gamma$ of $X_H$ gives the integral curve $\lambda\gamma$ of $X_{\tilde H}$, and
$\varphi^t_{\tilde H}=\delta_\lambda\circ\varphi^t_H\circ\delta_\lambda^{-1}$ wherever either side is defined.
\end{lemma}

\begin{proof}
The first two statements are immediate. Put $X(t,z)=\lambda X_H(t,z/\lambda)$. For $w\in\C^n$,
$\omega_0(X(t,z),w)=\lambda\,dH_{(t,z/\lambda)}(w)$, while $d\tilde H_{(t,z)}(w)=\lambda^2\,dH_{(t,z/\lambda)}(w/\lambda)$ is the same
number; so $X=X_{\tilde H}$. If $\dot\gamma=X_H(t,\gamma)$ then
$\frac d{dt}(\lambda\gamma)=\lambda X_H(t,\gamma)=X_{\tilde H}(t,\lambda\gamma)$.
\end{proof}

With $\lambda=1/\sqrt C$ this is the normalisation of the target capacity to $1$ used throughout
(\Msec{sec:ham}).

\emph{Moments and angles.} On a fibre $\C^m$ with coordinates $w_j=\xi_j+i\eta_j$ we use
\[
 p_j=\pi|w_j|^2,\qquad \theta_j=\frac{\arg w_j}{2\pi}\in\R/\Z,\qquad |p|=\sum_{j=1}^mp_j,\qquad
 \omega_V=\sum_{j=1}^mdp_j\wedge d\theta_j=\sum_{j=1}^md\xi_j\wedge d\eta_j .
\]
The identity $dp_j\wedge d\theta_j=d\xi_j\wedge d\eta_j$ holds where $w_j\ne0$: in polar coordinates
$w_j=re^{i\phi}$ both sides equal $r\,dr\wedge d\phi$. The one-form $p_j\,d\theta_j=\frac12(\xi_j\,d\eta_j-\eta_j\,d\xi_j)$
is smooth on all of $\C$, with differential $d\xi_j\wedge d\eta_j$. A function $f(p)$ with $f$ smooth on an
open subset of $\R^m$ is a smooth function of $w$, because each $p_j$ is a polynomial in $(\xi_j,\eta_j)$;
such functions are called \emph{toric}.

\emph{Fibre Hamiltonians.} On fibres, and on the models built from them, Hamiltonian fields of fibre
Hamiltonians $Q$ follow the convention of \cite{OpenAI2026},
\[
 \iota_{X_Q}\omega_V=-dQ ,
\]
which is opposite to the convention for $H$. With this convention $X_{p_j}=\partial_{\theta_j}$ generates
$w_j\mapsto e^{2\pi it}w_j$, the rotation of period one in the positive direction, because
$\iota_{\partial_{\theta_j}}(dp_j\wedge d\theta_j)=-dp_j$. The field $X_Q$ of the fibre convention is the field
$X_{-Q}$ of the convention of $\omega_0$. Two functions \emph{Poisson-commute} if $\omega_V(X_F,X_G)=0$; this
condition does not depend on the convention, and it holds if and only if $F$ is constant along the
flow of $X_G$. The convention $\iota_{X_Q}\omega_V=-dQ$ is used for fibre Hamiltonians, for Hamiltonians
on the toric models and for Moser fields; the convention $\iota_{X_H}\omega_0=dH$ is used for the global Hamiltonian
on $\R^{2n}$ (\Msec{sec:ham}, \Msec{sec:eff}, \S\ref{sec:iso}, \S\ref{sec:cpl}).

\emph{Complex conventions and the Fubini--Study form.} We use
\[
 d^c=\frac{i}{4\pi}(\bar\partial-\partial),\qquad dd^c=\frac{i}{2\pi}\partial\bar\partial ,
\]
so that $d^cu$ is a real one-form for real $u$. Since $\partial$ and $\bar\partial$ are derivations,
$d^cF(u)=F'(u)\,d^cu$ for smooth $F$ and smooth real $u$. For a function $u=F(|z|^2)$ of one complex variable,
$\bar\partial u-\partial u=F'(|z|^2)(z\,d\bar z-\bar z\,dz)=-2i\,F'(|z|^2)|z|^2\,d\phi$ with $\phi=\arg z$, hence
\begin{equation}\label{eq:pre-dc}
 d^cF(|z|^2)=F'(|z|^2)\,|z|^2\,d\theta,\qquad\text{in particular}\qquad d^c\log|z|^2=d\theta,
\end{equation}
with $\theta=\arg z/2\pi$. The one-form $|z|^2d\theta=(x\,dy-y\,dx)/2\pi$ is smooth, so \eqref{eq:pre-dc} holds
on all of $\C$ when $F$ is smooth, and on $\C\setminus0$ for $F=\log$. Since $\int_{|z|=\varepsilon}d\theta=1$,
the current $dd^c\log|z|^2$ has unit mass at the origin. On $\CP^n$ we put
$\omega_{\rm FS}=dd^c\log\|Z\|^2$.

\begin{lemma}[The affine Fubini--Study chart]\label{lem:pre-fs}
\lean{OAI.PackingSufficiencySupport.Hamiltonian.affineFSPrimitive_smooth,
OAI.PackingSufficiencySupport.Hamiltonian.affineFSPrimitive_exterior,
OAI.PackingSufficiencySupport.Hamiltonian.fsBallMap_smooth,
OAI.PackingSufficiencySupport.Hamiltonian.fsBallMap_pullback,
OAI.PackingSufficiencySupport.Hamiltonian.fsBallMap_image,
OAI.PackingSufficiencySupport.Hamiltonian.fsAffineMap_smoothOn,
OAI.PackingSufficiencySupport.Hamiltonian.fsAffineMap_fsBallMap,
OAI.PackingSufficiencySupport.Hamiltonian.fsBallMap_fsAffineMap,
OAI.PackingSufficiencySupport.Hamiltonian.fsBallMap_primitive,
OAI.PackingSufficiencySupport.Hamiltonian.fsBallMap_gap}
On $\C^N$, put $D=1+|z|^2$ and
\[
 \lambda_{\rm FS}=
 \frac{1}{2\pi D}\sum_{j=1}^N(x_j\,dy_j-y_j\,dx_j),\qquad
 \omega_{\rm FS}=d\lambda_{\rm FS}.
\]
Then:
\begin{enumerate}[label=(\alph*)]
\item $\lambda_{\rm FS}=d^c\log D$ is smooth. Where every coordinate is nonzero it equals
 $\sum_j p_j^{\rm FS}\,d\theta_j$, with $p_j^{\rm FS}=|z_j|^2/D$ and
 $\theta_j=\arg z_j/(2\pi)$; its differential is $\sum_jdp_j^{\rm FS}\wedge d\theta_j$ there.
\item The map $\mathrm{ch}(z)=z/\sqrt{\pi D}$ is a smooth bijection of $\C^N$ onto $\Bo(1)$ with $\mathrm{ch}^*\omega_0=\omega_{\rm FS}$,
 with smooth inverse $w\mapsto\sqrt\pi\,w/\sqrt{1-\pi|w|^2}$.
\item $\mathrm{ch}^*\lambda_{\rm std}=\lambda_{\rm FS}$, where
 $\lambda_{\rm std}=\frac12\sum_j(x_j\,dy_j-y_j\,dx_j)$, and
 $\pi|\mathrm{ch}(z)|^2=1-D^{-1}<1$.
\end{enumerate}
\end{lemma}

\begin{proof}
The identity $d^c|z|^2=\sum_j(x_j\,dy_j-y_j\,dx_j)/(2\pi)$ and the chain rule give (a).
The inverse in (b) follows by solving $\pi|w|^2=|z|^2/(1+|z|^2)$ for the radial factor.
Both maps are smooth on their stated domains. Under a variable real radial rescaling $z\mapsto r(z)z$,
$\lambda_{\rm std}$ pulls back to $r(z)^2\lambda_{\rm std}$: the derivative of $r$ contributes
$\omega_0(z,z)=0$. Taking $r=(\pi D)^{-1/2}$ gives (c), and applying $d$ gives $\mathrm{ch}^*\omega_0=\omega_{\rm FS}$.
\end{proof}

The model forms of \Mref{Table}{tab:mech-models} are written with $d^c$ and $\omega_{\rm FS}$ in this
normalisation.

\emph{Bidegree and orientation.} Let $U$ be an open subset of a surface $\Sigma$ with oriented coordinates
$(x,y)$ over which a bundle $E=L_1\oplus\dots\oplus L_m$ has a unitary frame, so that $E|_U=U\times\C^m$. Every
form on an open subset of $U\times\C^m$ splits uniquely into components of bidegree $(r,s)$, where $r$
counts $dx,dy$ and $s$ counts $d\xi_j,d\eta_j$; the coefficients may depend on all variables. A form of
bidegree $(r,0)$ is called \emph{basic}, and a component vanishes if $r>2$ or $s>2m$. A top-degree form is
\emph{positive} if it is a positive multiple of $\omega_V^m\wedge dx\wedge dy$ for an oriented coordinate system
$(x,y)$ of $\Sigma$; when $\kappa$ is a positive area form on $\Sigma$, this is the orientation of
$\omega_V^m\wedge\kappa$. A one-form on $E$ is \emph{horizontal} if it annihilates the vertical tangent vectors.

\subsection{The flat step}\label{sec:pre-flat}

We use the following flat step for clocks and cutoffs, composing it with the explicitly defined
smooth functions and taking products.

\begin{definition}[Flat step]\label{def:pre-flat}
\lean{expNegInvGlue, Real.smoothTransition, SymplecticPacking.flatStep}
Let $a_\flat(x)=0$ for $x\le0$ and $a_\flat(x)=e^{-1/x}$ for $x>0$, and
\[
 S(x)=\frac{a_\flat(x)}{a_\flat(x)+a_\flat(1-x)}\qquad(x\in\R).
\]
\end{definition}

The denominator is positive because $x>0$ or $1-x>0$. On $(0,1)$,
$S(x)=e^{-1/x}/(e^{-1/x}+e^{-1/(1-x)})$; $S=0$ on $(-\infty,0]$ and $S=1$ on $[1,\infty)$.

\begin{lemma}[Properties of $S$]\label{lem:pre-flat}
\lean{Real.smoothTransition.contDiff,
Real.smoothTransition.zero_of_nonpos,
Real.smoothTransition.one_of_one_le,
Real.smoothTransition.monotone,
Real.smoothTransition.nonneg,
Real.smoothTransition.le_one,
OAI.PackingSufficiencySupport.Hamiltonian.smoothTransition_deriv_pos,
OAI.PackingSufficiencySupport.Hamiltonian.smoothTransition_deriv_eq_zero}
\uses{def:pre-flat}
$S$ is $C^\infty$ on $\R$, takes values in $[0,1]$, is nondecreasing on $\R$ and strictly increasing on
$[0,1]$, and $S(1-x)=1-S(x)$. For every $j\ge1$, $S^{(j)}=0$ on $(-\infty,0]\cup[1,\infty)$; in particular
$S^{(j)}(0)=S^{(j)}(1)=0$.
\end{lemma}

\begin{proof}
On $x>0$ every derivative of $a_\flat$ has the form $Q_j(1/x)e^{-1/x}$ with $Q_j$ a polynomial, and
$Q_j(1/x)e^{-1/x}\to0$ as $x\downarrow0$; by induction all derivatives of $a_\flat$ exist at $0$ and vanish there,
so $a_\flat$ is $C^\infty$. Since the denominator of $S$ is smooth and positive, $S$ is $C^\infty$. The identity
$S(1-x)=a_\flat(1-x)/(a_\flat(1-x)+a_\flat(x))=1-S(x)$ is immediate. On $(0,1)$,
$S(x)=1/(1+\exp(1/x-1/(1-x)))$ and the exponent is strictly decreasing, so $S$ is strictly increasing on
$[0,1]$; with the constant values outside $(0,1)$ this gives monotonicity and $0\le S\le1$. Finally $S$ is
constant on $(-\infty,0]$ and on $[1,\infty)$, so $S^{(j)}$, $j\ge1$, vanishes on the interiors of these sets, and
by continuity of $S^{(j)}$ also at $0$ and $1$.
\end{proof}

The vanishing of all derivatives at $0$ and $1$ is what allows paths to be concatenated in time without
corners (Lemma~\ref{lem:pre-concat}).

\begin{lemma}[Flat concatenation]\label{lem:pre-concat}
\uses{lem:pre-flat}
Let $0=t_0<t_1<\dots<t_N=1$ and let $\sigma_j(t)=S\big((t-t_{j-1})/(t_j-t_{j-1})\big)$, $1\le j\le N$, be the
\emph{clocks}.
\begin{enumerate}[label=(\alph*)]
\item Let $\mathcal P$ and $M$ be smooth manifolds and $\gamma_1,\dots,\gamma_N\colon[0,1]\times\mathcal P\to M$ smooth maps
(smooth up to the boundary, i.e.\ restrictions of smooth maps on $(-\delta,1+\delta)\times\mathcal P$ for some $\delta>0$)
such that $\gamma_j(1,y)=\gamma_{j+1}(0,y)$ for all $y$. Define $\bar\gamma\colon\R\times\mathcal P\to M$ by
$\bar\gamma(t,y)=\gamma_j(\sigma_j(t),y)$ for $t\in[t_{j-1},t_j]$, $\bar\gamma(t,y)=\gamma_1(0,y)$ for $t\le0$ and
$\bar\gamma(t,y)=\gamma_N(1,y)$ for $t\ge1$. Then $\bar\gamma$ is well defined and $C^\infty$, and at every point
$(t_j,y)$, $0\le j\le N$, every partial derivative of $\bar\gamma$ (in charts) that contains at least one derivative in
$t$ vanishes.
\item Let $\psi^{(1)},\dots,\psi^{(N)}\colon[0,1]\times M\to M$ be smooth isotopies of $M$ (smooth up to the boundary as in
(a)) with $\psi^{(j)}_0=\mathrm{id}$, generated by smooth time-dependent vector fields $Y^{(j)}_s$, so that
$\partial_s\psi^{(j)}_s=Y^{(j)}_s\circ\psi^{(j)}_s$. Put $\Phi_0=\mathrm{id}$, $\Phi_j=\psi^{(j)}_1\circ\Phi_{j-1}$, and
\[
 \varphi_t=\psi^{(j)}_{\sigma_j(t)}\circ\Phi_{j-1}\qquad(t\in[t_{j-1},t_j]),
\]
$\varphi_t=\mathrm{id}$ for $t\le0$ and $\varphi_t=\Phi_N$ for $t\ge1$. Then $\varphi$ is a smooth isotopy with
$\varphi_0=\mathrm{id}$ and $\varphi_1=\Phi_N=\psi^{(N)}_1\circ\dots\circ\psi^{(1)}_1$, generated by the field
$Y_t=\sigma_j'(t)\,Y^{(j)}_{\sigma_j(t)}$ on $[t_{j-1},t_j]$ and $Y_t=0$ for $t\notin(0,1)$. The field $Y$ is smooth on
$\R\times M$, and $Y$ and all its partial derivatives vanish at $t=t_j$. If $M$ is symplectic and
$Y^{(j)}_s$ is the Hamiltonian field of $Q^{(j)}_s$ (in either sign convention), with $Q^{(j)}$ smooth on $[0,1]\times M$
(up to the boundary), then $Y_t$ is the Hamiltonian field of $Q_t=\sigma_j'(t)\,Q^{(j)}_{\sigma_j(t)}$ on $[t_{j-1},t_j]$,
$Q_t=0$ for $t\notin(0,1)$, and $Q$ is smooth on $\R\times M$.
\end{enumerate}
\end{lemma}

\begin{proof}
(a) At $t=t_j$ the two definitions give $\gamma_j(1,y)=\gamma_{j+1}(0,y)$, and at $t=0,1$ the constant extensions agree
with $\gamma_1(0,y)$ and $\gamma_N(1,y)$; so $\bar\gamma$ is well defined and continuous. Fix $j$ and a point $(t_j,y_0)$, and
work in charts around $\bar\gamma(t_j,y_0)$ and $y_0$ (by continuity $\bar\gamma$ maps a neighbourhood into the chart). On
$[t_{j-1},t_j]\times\mathcal P$ the map $\bar\gamma$ is the smooth map $\gamma_j\circ(\sigma_j\times\mathrm{id})$. By the chain rule and
Fa\`a di Bruno's formula, a partial derivative of it containing at least one $t$-derivative is a sum of terms each of
which contains a factor $\sigma_j^{(i)}(t)$ with $i\ge1$; at $t=t_j$ (and at $t=t_{j-1}$) these factors are
$(t_j-t_{j-1})^{-i}S^{(i)}(1)=0$ (resp.\ $S^{(i)}(0)=0$) by Lemma~\ref{lem:pre-flat}. The derivatives containing only
$y$-derivatives are the $y$-derivatives of $\gamma_j(1,\cdot)=\gamma_{j+1}(0,\cdot)$ from both sides. Hence the one-sided
jets at $t=t_j$ from the left and from the right agree (the same holds at $t=0$ and $t=1$ against the constant
extensions). A continuous map which is $C^\infty$ up to the boundary on two closed half-spaces and whose
one-sided jets agree along the common boundary is $C^\infty$; so $\bar\gamma$ is $C^\infty$.

(b) The maps $\gamma_j(s,x)=\psi^{(j)}_s(\Phi_{j-1}(x))$ satisfy the hypotheses of (a), since
$\gamma_j(1,\cdot)=\Phi_j=\psi^{(j+1)}_0\circ\Phi_j=\gamma_{j+1}(0,\cdot)$; so $\varphi$ is smooth, and each $\varphi_t$ is a
diffeomorphism as a composition of diffeomorphisms. On $[t_{j-1},t_j]$,
\[
 \partial_t\varphi_t=\sigma_j'(t)\,(\partial_s\psi^{(j)})_{\sigma_j(t)}\circ\Phi_{j-1}
 =\sigma_j'(t)\,Y^{(j)}_{\sigma_j(t)}\circ\psi^{(j)}_{\sigma_j(t)}\circ\Phi_{j-1}=\sigma_j'(t)\,Y^{(j)}_{\sigma_j(t)}\circ\varphi_t .
\]
The field $(t,x)\mapsto\sigma_j'(t)Y^{(j)}(\sigma_j(t),x)$ is smooth on $[t_{j-1},t_j]\times M$, and every partial derivative of
it (in $t$, in $x$, or mixed) at $t=t_{j-1}$ or $t=t_j$ is a sum of terms containing a factor $\sigma_j^{(i)}$ with $i\ge1$,
hence vanishes there; on $t\notin(0,1)$ the field is $0$. So the one-sided jets at the seams are all zero and
$Y$ is smooth. The statement on Hamiltonians follows because $X_{cQ}=cX_Q$ for constants $c$ in either
convention, and $Q$ is smooth by the same argument applied to the smooth function $Q^{(j)}$ on $[0,1]\times M$.
\end{proof}

\subsection{Cutoffs}\label{sec:pre-cutoff}

\begin{definition}[Radial and interval cutoffs]\label{def:pre-cutoff}
\lean{SymplecticPacking.General.Ham.radialCutoff}
For $0\le r<r'$ and $u$ in a Euclidean space $\R^d$ put
\[
 \chi_{r,r'}(u)=1-S\Big(\frac{|u|^2-r^2}{r'^2-r^2}\Big).
\]
For $\ell_0<\ell_1$ and $0<w<(\ell_1-\ell_0)/2$ put
$\chi^{\rm I}_{\ell_0,\ell_1,w}(x)=S\big((x-\ell_0)/w\big)\,S\big((\ell_1-x)/w\big)$ for $x\in\R$.
\end{definition}

The argument of $S$ in $\chi_{r,r'}$ is a polynomial in $u$, so every derivative of a cutoff is an
explicit expression in the derivatives of $S$. We write
\[
 s_1=\max_{\R}S' ,
\]
a finite constant determined by $S$, since $S'$ is continuous and vanishes outside $[0,1]$.

\begin{lemma}[Cutoffs]\label{lem:pre-cutoff}
\lean{SymplecticPacking.General.Ham.radialCutoff_contDiff,
SymplecticPacking.General.Ham.radialCutoff_eq_one,
SymplecticPacking.General.Ham.radialCutoff_eq_zero,
SymplecticPacking.General.Ham.radialCutoff_mem_Icc,
SymplecticPacking.General.Ham.tsupport_radialCutoff_subset}
\uses{def:pre-cutoff, lem:pre-flat}
\begin{enumerate}[label=(\alph*)]
\item $\chi_{r,r'}$ is $C^\infty$, $O(d)$-invariant and nonincreasing in $|u|$, takes values in $[0,1]$, equals $1$
on $\EB(r)$ and $0$ outside $\EB^\circ(r')$, and all its derivatives of order $\ge1$ vanish on
$\EB(r)\cup(\R^d\setminus\EB^\circ(r'))$, in particular on the spheres $|u|=r$ and $|u|=r'$. Its gradient is
\[
 \nabla\chi_{r,r'}(u)=-\frac{2u}{r'^2-r^2}\,S'\Big(\frac{|u|^2-r^2}{r'^2-r^2}\Big),\qquad
 |\nabla\chi_{r,r'}|\le\frac{2r's_1}{r'^2-r^2}.
\]
\item $\chi^{\rm I}_{\ell_0,\ell_1,w}$ is $C^\infty$ with values in $[0,1]$, equals $1$ on $[\ell_0+w,\ell_1-w]$ and $0$ outside
$(\ell_0,\ell_1)$, and all its derivatives of order $\ge1$ vanish outside $(\ell_0,\ell_0+w)\cup(\ell_1-w,\ell_1)$.
\end{enumerate}
\end{lemma}

\begin{proof}
(a) Write $\chi_{r,r'}=1-S\circ\ell$ with $\ell(u)=(|u|^2-r^2)/(r'^2-r^2)$, a polynomial. Smoothness, invariance,
monotonicity and the range follow from Lemma~\ref{lem:pre-flat}. If $|u|\le r$ then $\ell(u)\le0$ and $\chi=1$; if
$|u|\ge r'$ then $\ell(u)\ge1$ and $\chi=0$. By Fa\`a di Bruno's formula every derivative of order $\ge1$ of
$S\circ\ell$ at $u$ is a sum of terms each containing a factor $S^{(i)}(\ell(u))$ with $i\ge1$, which vanishes when
$\ell(u)\le0$ or $\ell(u)\ge1$ (Lemma~\ref{lem:pre-flat}). The gradient is the chain rule, and $|u|\le r'$ wherever
$S'(\ell(u))\ne0$.
(b) If $x\le\ell_0$ or $x\ge\ell_1$ one factor vanishes. If $\ell_0+w\le x\le\ell_1-w$ both arguments are $\ge1$. On
$(\ell_0,\ell_0+w)$ we have $x<\ell_1-w$ because $2w<\ell_1-\ell_0$, so the second factor is $1$ and only the first varies;
symmetrically on $(\ell_1-w,\ell_1)$. The statement on derivatives follows as in (a).
\end{proof}

A cutoff around a coordinate disc always means the squared-distance cutoff $\chi_{r,r'}$ in that coordinate.

\subsection{Toric forms}\label{sec:pre-toric}

\begin{lemma}[Top power]\label{lem:pre-top-power}
\lean{OAI.PackingSufficiencySupport.Hamiltonian.toric_twoCovectorForm_isInvertible,
OAI.PackingSufficiencySupport.Hamiltonian.twoCovectorForm_isInvertible,
OAI.PackingSufficiencySupport.Hamiltonian.toric_horizontalCoupling_isInvertible,
OAI.PackingSufficiencySupport.Hamiltonian.oneCovectorForm_isInvertible,
SymplecticPacking.Forms.pf_splitForm}
\uses{M-def:hdl-handle-data}
Let $m\ge1$, $n=m+1$, and let $U$ be an oriented surface chart with coordinates $(x,y)$.
On an open subset of $U\times\C^m$, write $w_j=\xi_j+i\eta_j$,
$\omega_V=\sum_jd\xi_j\wedge d\eta_j$, and $\kappa'=\varkappa\,dx\wedge dy$.
The mixed term $\vartheta$ is of one of the two types used in the construction:
\begin{enumerate}[label=(\alph*)]
\item $\vartheta=\sum_jdp_j\wedge a'_j$, where $p_j=\pi|w_j|^2$ and the $a'_j$ are basic one-forms;
 their coefficients may depend on the base and fibre variables.
\item $\vartheta=\zeta\wedge\beta'$, where $\zeta$ is vertical and $\beta'$ is basic.
\end{enumerate}
In both cases $\omega_V^{m-1}\wedge\vartheta^2=0$, and
\[
 (\omega_V+\kappa'+\vartheta)^n
 =n\,\omega_V^m\wedge\kappa'
 =n!\,\varkappa\;d\xi_1\wedge d\eta_1\wedge\dots\wedge d\xi_m\wedge d\eta_m\wedge dx\wedge dy.
\]
Thus the two-form is nondegenerate exactly where $\varkappa\ne0$, and positive where $\varkappa>0$.
In particular, on the handle annulus take $\beta=b(v)\,dv$ closed, $\nu=c_\nu\,ds\wedge dv$,
$e>0$, $c_\nu>0$, $b\ge0$, and a smooth $G(s,w)$ independent of $v$. Then
\[
 \big(\omega_V+e\nu+d(G\beta)\big)^n
 =n\big(e c_\nu+b(v)\partial_sG\big)\omega_V^m\wedge ds\wedge dv.
\]
It is symplectic and positive where $\partial_sG\ge0$.
\Mref{Lemma}{lem:hdl-top-power} uses precisely these two cases.
\end{lemma}

\begin{proof}
Forms of even degree commute, so the multinomial formula gives
\[
 (\omega_V+\kappa'+\vartheta)^n=\sum_{i_1+i_2+i_3=n}\frac{n!}{i_1!\,i_2!\,i_3!}\,\omega_V^{i_1}\wedge\kappa'^{i_2}\wedge\vartheta^{i_3} ,
\]
and the summand of index $(i_1,i_2,i_3)$ has bidegree $(2i_2+i_3,\,2i_1+i_3)$. It vanishes unless $2i_2+i_3\le2$ and
$2i_1+i_3\le2m$. The first condition leaves $(i_2,i_3)\in\{(0,0),(0,1),(0,2),(1,0)\}$, and then $i_1=n-i_2-i_3$. For
$(i_2,i_3)=(0,0)$ and $(0,1)$ the fibre degree is $2n$ and $2m+1$, both larger than $2m$. For $(i_2,i_3)=(1,0)$, $i_1=m$,
and the summand is $n\,\omega_V^m\wedge\kappa'$. For $(i_2,i_3)=(0,2)$, $i_1=m-1$, and the summand is
$\frac{n!}{2\,(m-1)!}\,\omega_V^{m-1}\wedge\vartheta^2$, which vanishes by the identity proved below in cases (a) and (b). Since
$\omega_V^m=m!\,d\xi_1\wedge d\eta_1\wedge\dots\wedge d\xi_m\wedge d\eta_m$, this proves the formula. A two-form on a vector space of
dimension $2n$ is nondegenerate if and only if its $n$-th power is nonzero.

(a) $\vartheta\wedge\vartheta=-\sum_{j,l}dp_j\wedge dp_l\wedge a'_j\wedge a'_l$, and
$\omega_V^{m-1}=(m-1)!\sum_{i=1}^m\bigwedge_{r\ne i}d\xi_r\wedge d\eta_r$. The one-form
$dp_j=2\pi(\xi_jd\xi_j+\eta_jd\eta_j)$ lies in the span of $d\xi_j,d\eta_j$, so $d\xi_j\wedge d\eta_j\wedge dp_j=0$. For $j=l$,
$dp_j\wedge dp_j=0$. For $j\ne l$ and any $i$, one of $j,l$ differs from $i$, say $j$; then the summand
$\bigwedge_{r\ne i}d\xi_r\wedge d\eta_r$ contains $d\xi_j\wedge d\eta_j$, and its product with $dp_j$ vanishes. Hence
the mixed-square identity. (For $m=1$ only $j=l=1$ occurs.)

(b) $(\zeta\wedge\beta')\wedge(\zeta\wedge\beta')=-\zeta\wedge\zeta\wedge\beta'\wedge\beta'=0$ because $\beta'\wedge\beta'=0$. For the displayed
case: $d\beta=0$, $\beta=b(v)\,dv$ and $\partial_vG=0$, so $d(G\beta)=dG\wedge\beta=\partial_sG\,b(v)\,ds\wedge dv+d_wG\wedge\beta$, where
$d_wG=\sum_j(\partial_{\xi_j}G\,d\xi_j+\partial_{\eta_j}G\,d\eta_j)$ has bidegree $(0,1)$. With $\nu=c_\nu\,ds\wedge dv$ the form is
$\omega_V+\kappa'+\vartheta$ with $\kappa'=(ec_\nu+b\,\partial_sG)\,ds\wedge dv$ and $\vartheta=d_wG\wedge\beta$, of the type just treated.
Since $e>0$, $c_\nu>0$ and $b\ge0$, the coefficient is positive where $\partial_sG\ge0$.
\end{proof}

The non-toric term $d_wG\wedge\beta$ of (b) is harmless because it is divisible by the single basic one-form
$\beta$, so it squares to zero (closedness of $\beta$ is used only for $d(G\beta)=dG\wedge\beta$). A sum of a toric mixed
term and a term $\zeta\wedge\beta$ need not satisfy the mixed-square identity: for $m=1$,
$\omega_V^0\wedge(dp\wedge a'+\zeta\wedge\beta)^2=2\,dp\wedge a'\wedge\zeta\wedge\beta$ is in general nonzero. The handle construction
therefore uses (a) off the annulus $\mathcal A$ and (b) on it, never a mixture.

The model calculations below are local coordinate calculations. On an oriented base chart $U$,
let $a_j$ be smooth real one-forms, $F_j=-da_j$, and let $\sigma$ be a smooth two-form.
On $U\times\C^m$ put $p_j=\pi|w_j|^2$ and
\[
 \lambda_j=\tfrac12(\xi_j\,d\eta_j-\eta_j\,d\xi_j),\qquad
 p_j\alpha_j=\lambda_j+p_ja_j.
\]
Here $\alpha_j=d\theta_j+a_j$ where $w_j\ne0$; the product $p_j\alpha_j$ is defined everywhere by
the displayed smooth expression. These formulas require no choice of a global frame or identification
with a projective normal bundle. In a cover presentation their descent is checked by the stated deck transformations.

\begin{lemma}[Local form of the model form]\label{lem:pre-local-form}
\uses{lem:pre-fs, lem:pre-top-power}
On the above product chart put
\[
 \Omega_0=\sigma+\sum_{j=1}^md(p_j\alpha_j).
\]
\begin{enumerate}[label=(\alph*)]
\item The one-forms $p_j\alpha_j$ are smooth, and $\Omega_0$ is closed and invariant under fibre rotations.
\item
\[
 \Omega_0=\omega_V+\kappa(p)+\sum_jdp_j\wedge a_j,
 \qquad \kappa(p)=\sigma-\sum_jp_jF_j.
\]
\item For the circle generator $X_{p_j}=\partial_{\theta_j}$,
 $\iota_{X_{p_j}}\Omega_0=-dp_j$, including the fibre axis.
\item With $n=m+1$, $\Omega_0^n=n\,\omega_V^m\wedge\kappa(p)$.
 Thus it is nondegenerate where $\kappa(p)\ne0$ and positive where $\kappa(p)>0$.
 If $F_j=d_j\sigma$ for constants $d_j$, its horizontal coefficient is
 $\kappa(p)=(1-\sum_jd_jp_j)\sigma$.
\item On a fixed compact base patch, integration at fixed moment gives
 $\int\kappa(p)=\int\sigma-\sum_jp_j\int F_j$.
\end{enumerate}
\end{lemma}

\begin{proof}
The defining expression $\lambda_j+p_ja_j$ is smooth, and
\[
 d(\lambda_j+p_ja_j)=d\xi_j\wedge d\eta_j+dp_j\wedge a_j-p_jF_j.
\]
This proves (a),(b), since two-forms on the two-dimensional base are closed.
Contraction with a fibre rotation gives (c); both sides vanish on its axis.
Part (d) is Lemma~\ref{lem:pre-top-power} for toric mixed terms.
Linearity of integration gives (e).
\end{proof}

For the periodic cubic input, $m=1$, $F_1=3\sigma$ and $\kappa=(1-3p)\sigma$; for the quadric curves,
$F_j=2\sigma$ and $\kappa=(1-2|p|)\sigma$ (\S\ref{sec:cub}, \S\ref{sec:qdr}).

\subsection{The Moser lemma}\label{sec:pre-moser}

\begin{lemma}[Moser lemma]\label{lem:pre-moser}
\lean{SymplecticPacking.Handle.LField.flow_preserves_form}
\uses{lem:pre-concat}
Let $M$ be a smooth manifold, $U\subset M$ open and $I\subset\R$ an open interval containing $[0,1]$. Let $\Omega_t$ and
$\mathfrak v_t$, $t\in I$, be two-forms and one-forms on $U$ depending smoothly on $(t,x)\in I\times U$, such that each
$\Omega_t$ is closed and nondegenerate at every point of $U$ and
\[
 \partial_t\Omega_t=d\mathfrak v_t .
\]
Let $Y_t$ be the vector field on $U$ defined by $\iota_{Y_t}\Omega_t=-\mathfrak v_t$, and let $\Psi_t$ be its flow from time $0$
($\Psi_0=\mathrm{id}$), where defined.
\begin{enumerate}[label=(\alph*)]
\item $Y$ is smooth on $I\times U$, and $Y_t(x)=0$ whenever $\mathfrak v_t(x)=0$ as an element of $T_x^*M$.
\item If $\Psi_t(x)$ is defined for $t\in[0,1]$, then $\Psi_t$ is defined and smooth on a neighbourhood of $x$ and
$(\Psi_t^*\Omega_t)_x=(\Omega_0)_x$ for $t\in[0,1]$.
\item If $Y_t(x)=0$ for all $t\in[0,1]$, then $\Psi_t(x)=x$ for $t\in[0,1]$. In particular $\Psi_t$ fixes pointwise every set
at whose points $\mathfrak v_t$ is the zero covector for all $t\in[0,1]$.
\item (Tangency.) Let $U'\subset U$ be open, $\mu\in C^\infty(U')$ and $V$ a vector field on $U'$ with
$\iota_V\Omega_t=-d\mu$ and $\mathfrak v_t(V)=0$ on $U'$ for $t\in[0,1]$. Then $d\mu(Y_t)=0$ on $U'$, and $\mu$ is constant along
every integral curve of $Y$ on a subinterval of $[0,1]$ as long as the curve stays in $U'$.
\item (Completeness.) Let $\mathcal K\subset U$ be compact, and suppose that every integral curve $\gamma\colon J\to U$ of $Y$,
$J\subset[0,1]$ an interval, which meets $\mathcal K$ satisfies $\gamma(J)\subset\mathcal K$. Then $\Psi_t$ is defined on $\mathcal K$ for
every $t\in[0,1]$, $\Psi_t(\mathcal K)=\mathcal K$, $\Psi_t|_{\mathcal K}$ is the restriction of a diffeomorphism between open
neighbourhoods of $\mathcal K$, and $\Psi_t^*\Omega_t=\Omega_0$ at every point of $\mathcal K$.
\item (Reparametrisation and concatenation.) If $\varsigma\colon I\to I$ is smooth with $\varsigma(0)=0$, the pair
$(\Omega_{\varsigma(t)},\varsigma'(t)\mathfrak v_{\varsigma(t)})$ satisfies the hypotheses, its field is $\varsigma'(t)Y_{\varsigma(t)}$, and its flow is
$\Psi_{\varsigma(t)}$. Consequently, let $(\Omega^{(j)}_s,\mathfrak v^{(j)}_s)$, $1\le j\le N$, be pairs on $U$, smooth on
$[0,1]\times U$ (up to the boundary), with every $\Omega^{(j)}_s$ closed and nondegenerate on $U$,
$\partial_s\Omega^{(j)}_s=d\mathfrak v^{(j)}_s$, and $\Omega^{(j)}_1=\Omega^{(j+1)}_0$. With the clocks $\sigma_j$ of
Lemma~\ref{lem:pre-concat} put $(\Omega_t,\mathfrak v_t)=(\Omega^{(j)}_{\sigma_j(t)},\sigma_j'(t)\mathfrak v^{(j)}_{\sigma_j(t)})$ on $[t_{j-1},t_j]$,
$(\Omega_t,\mathfrak v_t)=(\Omega^{(1)}_0,0)$ for $t\le0$ and $(\Omega_t,\mathfrak v_t)=(\Omega^{(N)}_1,0)$ for $t\ge1$. Then
$(\Omega_t,\mathfrak v_t)_{t\in\R}$ satisfies the hypotheses with $I=\R$, and its time-one map is
$\Psi_1=\Psi^{(N)}_1\circ\dots\circ\Psi^{(1)}_1$ wherever the right side is defined.
\item (Equivariance.) If a diffeomorphism $g$ of $U$ satisfies $g^*\Omega_t=\Omega_t$ and $g^*\mathfrak v_t=\mathfrak v_t$ for all $t$,
then $g_*Y_t=Y_t$ and $g\circ\Psi_t=\Psi_t\circ g$ wherever defined.
\end{enumerate}
\end{lemma}

\begin{proof}
(a) In local coordinates $\Omega_t$ is a smooth family of invertible skew matrices and $Y_t=-\Omega_t^{-1}\mathfrak v_t$; the
inverse is smooth by Cramer's rule, and $Y_t(x)=0$ if $\mathfrak v_t(x)=0$.

(b) The domain of the flow is open, so if $\Psi_t(x)$ is defined for $t\in[0,1]$, then $\Psi_t$ is defined and smooth on a
neighbourhood of $x$ for $t$ in a neighbourhood of $[0,1]$. For time-dependent flows,
$\frac d{dt}\Psi_t^*\Omega_t=\Psi_t^*\big(\partial_t\Omega_t+\mathcal L_{Y_t}\Omega_t\big)$. By Cartan's formula and $d\Omega_t=0$,
$\mathcal L_{Y_t}\Omega_t=d\iota_{Y_t}\Omega_t=-d\mathfrak v_t$, so the derivative vanishes, and $\Psi_0^*\Omega_0=\Omega_0$.

(c) The constant curve at $x$ is an integral curve; by uniqueness it is the integral curve through $x$. The second
statement follows by (a).

(d) $d\mu(Y_t)=-(\iota_V\Omega_t)(Y_t)=-\Omega_t(V,Y_t)=\Omega_t(Y_t,V)=(\iota_{Y_t}\Omega_t)(V)=-\mathfrak v_t(V)=0$, and
$\frac d{dt}\mu(\gamma(t))=d\mu(Y_t)=0$ along integral curves in $U'$.

(e) Choose a smooth compactly supported cutoff equal to $1$ near $[0,1]\times\mathcal K$, with support in
$I\times U$, and extend the cut-off field by zero. Its flow $\Phi_t$ and inverse are defined on all of $M$
for $t\in[0,1]$ and are jointly smooth near that interval.
For any trajectory of this field, the times at which it lies in $\mathcal K$ form a relatively closed subset
of $[0,1]$. This subset is also relatively open: near a point of $[0,1]\times\mathcal K$ the cutoff is $1$,
so the trajectory solves the original equation, and the trapping hypothesis applies on a time interval
about that point. Connectedness of $[0,1]$ shows that a trajectory meeting $\mathcal K$ stays there.
This applies both to a trajectory starting in $\mathcal K$ and to the trajectory through a point of
$\mathcal K$ at time $t$, proving $\Phi_t(\mathcal K)=\mathcal K$ in both directions.
The tube lemma gives an open neighbourhood of $\mathcal K$ whose trajectories remain where the cutoff
is $1$ throughout $[0,1]$. There $\Phi$ is the original flow $\Psi$, with its smooth inverse;
part (b) gives the pullback identity.

(f) $\iota_{\varsigma'(t)Y_{\varsigma(t)}}\Omega_{\varsigma(t)}=-\varsigma'(t)\mathfrak v_{\varsigma(t)}$ and $\partial_t\Omega_{\varsigma(t)}=\varsigma'(t)\,d\mathfrak v_{\varsigma(t)}$.
The curve $t\mapsto\Psi_{\varsigma(t)}(x)$ starts at $x$ and has derivative $\varsigma'(t)Y_{\varsigma(t)}(\Psi_{\varsigma(t)}(x))$, so by
uniqueness it is the flow of the reparametrised field. For the concatenation, each $\Omega_t$ is one of the closed
nondegenerate forms $\Omega^{(j)}_s$; $\Omega$ is smooth on $\R\times U$ by Lemma~\ref{lem:pre-concat}(a), applied to the maps
$(s,x)\mapsto\Omega^{(j)}_s(x)$ into the bundle of two-forms; and $\mathfrak v$ is smooth on $\R\times U$ by the argument of
Lemma~\ref{lem:pre-concat}(b): on each $[t_{j-1},t_j]$ every partial derivative of $\sigma_j'(t)\mathfrak v^{(j)}_{\sigma_j(t)}$, including
the zeroth, contains a factor $\sigma_j^{(i)}$ with $i\ge1$, so all one-sided $t$-jets vanish at every $t_j$, as they do for
the constant extension $0$. On $[t_{j-1},t_j]$, $\partial_t\Omega_t=\sigma_j'(t)\,(\partial_s\Omega^{(j)})_{\sigma_j(t)}=d\mathfrak v_t$, and off
$(0,1)$ both sides vanish. On $[t_{j-1},t_j]$ the curve $t\mapsto\Psi^{(j)}_{\sigma_j(t)}(x)$ starts at $x$ for $t=t_{j-1}$ and has derivative
$\sigma_j'(t)Y^{(j)}_{\sigma_j(t)}=Y_t$ at its points, as in the first part; so by uniqueness the flow of $Y$ from time $t_{j-1}$
to time $t$ is $\Psi^{(j)}_{\sigma_j(t)}$, and composing the stages gives $\Psi_1$.

(g) Since $g^*\mathfrak v_t=\mathfrak v_t$, also $(g^{-1})^*\mathfrak v_t=\mathfrak v_t$, and
$\iota_{g_*Y_t}\Omega_t=(g^{-1})^*\big(\iota_{Y_t}g^*\Omega_t\big)=(g^{-1})^*\big(\iota_{Y_t}\Omega_t\big)=-(g^{-1})^*\mathfrak v_t=-\mathfrak v_t$;
by uniqueness $g_*Y_t=Y_t$. Flows of $g$-related fields are conjugated by $g$.
\end{proof}

\emph{Signs.} The Moser field is defined by $\iota_{Y_t}\Omega_t=-\mathfrak v_t$, the same sign as the fibre convention: if
$\Omega_t=\omega$ is constant and $\mathfrak v_t=dQ_t$, then $Y_t=X_{Q_t}$ with $\iota_{X_Q}\omega=-dQ$. The field depends on the
chosen primitive $\mathfrak v_t$, not only on $\partial_t\Omega_t$; properties (c) and (d) are properties of the primitive.
Part (c) needs $\mathfrak v_t$ to vanish as a covector on $M$; vanishing of its pullback to a submanifold only makes
$Y_t$ $\Omega_t$-orthogonal to that submanifold.

\subsection{Conventions across the construction}\label{sec:pre-unified}

Table~\ref{tab:pre-conventions} lists the conventions by setting. Two of them differ in a way that matters for
checking: the global Hamiltonian uses $\iota_{X_H}\omega_0=dH$, while fibre Hamiltonians and Moser fields use the
opposite sign.
Table~\ref{tab:pre-interfaces} lists, for each interface between parts of the construction, the convention on each
side and the place where the conversion is made.

{\footnotesize
\renewcommand{\arraystretch}{1.12}
\begin{longtable}{@{}p{0.15\textwidth}p{0.62\textwidth}p{0.17\textwidth}@{}}
\caption{Conventions by setting.}\label{tab:pre-conventions}\\
\toprule
setting & convention & where\\\midrule
\endfirsthead
\multicolumn{3}{c}{\tablename\ \thetable\ (continued)}\\
\toprule
setting & convention & where\\\midrule
\endhead
\midrule
\multicolumn{3}{r}{Continued on the next page}\\
\endfoot
\bottomrule
\endlastfoot
phase space $\R^{2n}$ & $z_j=x_j+iy_j$; real coordinates $\zeta=(x_1,\dots,x_n,y_1,\dots,y_n)$, the order in which
 $J_0=\bigl(\begin{smallmatrix}0&I\\-I&0\end{smallmatrix}\bigr)$; $\omega_0=\sum_jdx_j\wedge dy_j$,
 $\omega_0(v,w)=\langle v,J_0w\rangle$; multiplication by $i$ is $J_\C=-J_0$ &
 \S\Mnum{sec:intro}, \S\Mnum{sec:ham}, \S\ref{sec:pre-conventions}, \S\ref{sec:iso}\\
global Hamiltonian & $\iota_{X_H}\omega_0=dH$, $X_H=J_0\nabla H$, $\dot x_j=\partial_{y_j}H$, $\dot y_j=-\partial_{x_j}H$;
 $H=\pi|z_1|^2$ rotates $z_1$ by $e^{-2\pi it}$. Used for $H[n,k,C,c,s]$, for $\Hfin=\beta_BP$ and for the
 certificate fields $J_0\nabla P_e$ & \S\Mnum{sec:intro}, \S\ref{sec:cpl}, \S\ref{sec:certificate}\\
capacities and radii & $B(c)=\{\pi|z|^2\le c\}$ has Euclidean radius $\sqrt{c/\pi}$ and volume $c^n/n!$;
 $B(s,c)=\EB(s,\sqrt{c/\pi})$. The models, transfers and comparison data work with capacities; the common
 isotopy works with the radii $r_i=\sqrt{c_i/\pi}$, $r_{\ell,i}$, $r'_i$, $\bar r_i$. In the Epoch problem the unit
 ball is $B(s_i,\pi)$ and $\Bo(\pi R^2)\subset\EB(R)$ &
 \S\Mnum{sec:intro}, Def.~\Mnum{def:ham-radii}, Cor.~\Mnum{thm:main-epoch}\\
normalisation & target capacity $1$: $\rho_i=c_i/C$; $\delta(z)=z/\sqrt C$ maps $\Bo(C)$ to $\Bo(1)$, admissibility is
 homogeneous, $\tilde H(t,z)=\lambda^2H(t,z/\lambda)$ under $z\mapsto\lambda z$ &
 Def.~\Mnum{def:ham-normalise}, Lemma~\ref{lem:pre-dilation}\\
toric fibre $\C^m$ & $w_j=\xi_j+i\eta_j$ (written $x_j+iy_j$ in \S\Mnum{sec:mech-toric}); $p_j=\pi|w_j|^2$,
 $\theta_j=\arg w_j/2\pi\in\R/\Z$, $|p|=\sum_jp_j$; $\omega_V=\sum_jd\xi_j\wedge d\eta_j=\sum_jdp_j\wedge d\theta_j$ &
 \S\Mnum{sec:intro-conventions}, \S\ref{sec:pre-conventions}\\
fibre Hamiltonian & $\iota_{X_Q}\omega_V=-dQ$; $X_{p_j}=\partial_{\theta_j}$ rotates $w_j$ by $e^{+2\pi it}$; the field
 $X_Q$ of this convention is $X_{-Q}$ of the global one. The same sign is used in the plane,
 $\iota_X\omega_\C=-dQ$, for the comparison data and the cubic data &
 \S\Mnum{sec:cmp}, Def.~\Mnum{def:mech-datum}, \S\ref{sec:pln}, Def.~\ref{def:shd4-data}\\
Moser fields & $\iota_{Y_t}\Omega_t=-\mathfrak v_t$ for $\partial_t\Omega_t=d\mathfrak v_t$; for constant $\Omega_t$ and
 $\mathfrak v_t=dQ_t$ this is the fibre convention. Used for the handle paths and for the returns of the transfers &
 Lemma~\ref{lem:pre-moser}\\
planar machine & disc $D(P)=\{\pi|z|^2<P\}$ of area $P$; $\omega_\C=dp\wedge d\vartheta$ with $p=\pi|z|^2$,
 $\vartheta=\arg z/2\pi$; Morse plane $(u,v)$ with $du\wedge dv$; rectangle $(x,\theta)$ with $dx\wedge d\theta$ and
 $\theta=\vartheta-\frac12$, chart $(v,\theta)\mapsto-\sqrt{v/\pi}\,e^{2\pi i\theta}$; symplectic means Jacobian $1$ &
 \S\ref{sec:pln}, \S\ref{sec:pll-setting}, \S\Mnum{sec:cmp-realisation}\\
model balls & $\C^n=\C\times\C^m\ni(z_{\rm b},z_{\rm f})$, $\varpi=\pi|z_{\rm b}|^2$, $B(c)=\{\varpi+|p|\le c\}$,
 $\omega_0=\omega_\C(z_{\rm b})+\omega_V(z_{\rm f})$; slice height $(c-|p|)_+$ &
 Def.~\Mnum{def:mech-toric-ball}, \S\Mnum{sec:hdl-proof}\\
model forms & $\Omega_0=\sigma+\sum_jd(p_j\alpha_j)=\omega_V+\kappa(p)+\sum_jdp_j\wedge a_j$, $\iota_{X_{p_j}}\Omega_0=-dp_j$;
 cubic cover $\R^2\times\C\ni(u,v,w)$: $\alpha=d\arg w/2\pi-3\eta$, $\Omega_0=\omega_V+(1-3p)\sigma-3\,dp\wedge\eta$,
 $\int_\Gamma\sigma=3$; quadric: $\Omega_0=\omega_V+d[(1-2|p|)\lambda]$; large ball, $n\ge4$: $u=p_1$,
 $w=\sum_{j\ge2}p_j$, $\Omega_{\rm mix}=\omega_V+d[(1-u-2w)\lambda]$ &
 Lemma~\ref{lem:pre-local-form}, Def.~\ref{def:cub-model}, Table~\Mnum{tab:mech-models}, \S\ref{sec:lrg-high}\\
K\"ahler forms & $d^c=\frac i{4\pi}(\bar\partial-\partial)$, so $dd^c\log|z|^2$ has unit mass and a line in $\CP^n$ has
 area $1$; $\omega_{\rm FS}=dd^c\log\|Z\|^2$, in the affine chart $dd^c\log(1+|z|^2)$;
 $\mathrm{ch}(z)=z/\sqrt{\pi(1+|z|^2)}$ maps $\C^N$ onto $\Bo(1)$ with $\mathrm{ch}^*\omega_0=\omega_{\rm FS}$; cubic
 representatives $(Z,W)\sim(e^{i\chi}Z,e^{3i\chi}W)$ &
 Lemma~\ref{lem:pre-fs}, \S\ref{sec:rtc-conventions}\\
\end{longtable}
}

{\footnotesize
\renewcommand{\arraystretch}{1.12}
\begin{longtable}{@{}p{0.17\textwidth}p{0.6\textwidth}p{0.17\textwidth}@{}}
\caption{Conventions at the interfaces and where they are converted.}\label{tab:pre-interfaces}\\
\toprule
interface & conventions and conversion & where\\\midrule
\endfirsthead
\multicolumn{3}{c}{\tablename\ \thetable\ (continued)}\\
\toprule
interface & conventions and conversion & where\\\midrule
\endhead
\midrule
\multicolumn{3}{r}{Continued on the next page}\\
\endfoot
\bottomrule
\endlastfoot
comparison $\to$ planar realisation & The comparison is a statement on the moment polytope, in the moments $p$;
 the rearrangement $R_j$ acts on $p_j$. The planar map lives on a disc of area $1$ in the rectangle chart, with the
 fibre sign on both sides. The conversion is the rescaling of the inner slice to area $1$:
 $Q^{(j)}_t(w)=\chi(y)S'(t)\ell_{\rm in}(y)\,Q^{\rm pl}\bigl(y,\chi(y)S(t),w_j/\sqrt{\ell_{\rm in}(y)}\bigr)$, the dilation
 rule with $\lambda^2=\ell_{\rm in}$. For $n=2$ no rescaling occurs: the planar path acts on the fibre $\C$ itself &
 \S\Mnum{sec:cmp-realisation}, Thm.~\ref{thm:plc-lift}, Prop.~\ref{prop:scalar-planar-realisation}\\
fibre Hamiltonian $\to$ bundle & The isotopy $q_t$, generated with the fibre sign, is suspended over the annulus by
 $\mathcal D_\lambda(s,v,w)=(s,v,q_{\lambda f(s)}(w))$ with generator $Q^\lambda_s=\lambda f'(s)Q_{\lambda f(s)}$; then
 $\mathcal D_\lambda^*\omega_V=\omega_V+d[(Q^\lambda_s\circ q_{\lambda f(s)})\,ds]$. The time of $q$ becomes the annulus
 coordinate $s$; the sign is unchanged, and the Moser fields of the three paths use the same sign &
 Lemma~\Mnum{lem:hdl-suspension}, \S\ref{sec:shd-four}\\
model $\to$ transfer to $\Bo(1)$ & The model embeddings satisfy $\psi_i^*\Omega=\omega_0$ for the splitting
 $\C\times\C^m$ of the source; the transfer satisfies $\mathcal T^*\omega_0=\Omega$ (resp.\ $\Omega_0$) and lands in
 $\Bo(1)$, the target normalised to capacity $1$. The passage from $\omega_{\rm FS}$ to $\omega_0$ is the chart
 $\mathrm{ch}$ (quadric) or $b_{\rm ball}$, with $b_{\rm ball}^*\omega_0=\omega_{\rm FS}$ on unit representatives (cubic) &
 Prop.~\Mnum{prop:mech-assembly}, Thms.~\ref{thm:cub-transfer}, \ref{thm:qdr-transfer}, \S\ref{sec:rtc-avoid}\\
embeddings $\to$ common isotopy & The models give normalised maps $f^\natural_i$ for $(n,k,1,\rho)$ with capacities
 $\rho_i<\rho_i^+<\rho_i^{\rm dom}$. The conversion to the physical target is
 $f_i(u)=\sqrt C\,f^\natural_i(u/\sqrt C)$, $c'_i=C\rho_i^+$, $\bar c_i=C\rho_i^{\rm dom}$; from here on balls are
 described by Euclidean radii and linear maps by matrices in $(x,y)$ blocks &
 Def.~\Mnum{def:ham-normalise}, Prop.~\Mnum{prop:ham-embeddings}, Def.~\Mnum{def:ham-radii}\\
Hamiltonian $\to$ finite form & Same coordinates and sign. $P\in\Q[t,z]$ is a polynomial in $t$ and the real
 coordinates $\zeta_1,\dots,\zeta_{2n}$ (the variables $x_0=t,x_1,\dots,x_{2n}$ of a polynomial program), and
 $\beta_B(z)=\prod_l\beta_{B,1}(\zeta_l)$. The Epoch interface reads $(Q,P)=(x,y)$ &
 \S\Mnum{sec:eff}, Def.~\ref{def:certificate-program}, Rem.~\Mnum{rem:eff-epoch-interface}\\
\end{longtable}
}

\subsection{Notation}\label{sec:pre-notation}

Table~\ref{tab:pre-notation} lists the symbols used in more than one section. Local symbols are introduced
where they are used; a statement that reuses a letter of the table redefines it there.

{\footnotesize
\renewcommand{\arraystretch}{1.12}
\begin{longtable}{@{}p{0.27\textwidth}p{0.55\textwidth}p{0.12\textwidth}@{}}
\caption{Notation used in more than one section. The letters $s$, $y$, $b$ and $J$ have more than one row
or a local meaning (for instance in Table~\Mnum{tab:mech-models} and in \S\ref{sec:lrg}); the section fixes which is meant.}
\label{tab:pre-notation}\\
\toprule
symbol & meaning & defined in\\\midrule
\endfirsthead
\multicolumn{3}{c}{\tablename\ \thetable\ (continued)}\\
\toprule
symbol & meaning & defined in\\\midrule
\endhead
\midrule
\multicolumn{3}{r}{Continued on the next page}\\
\endfoot
\bottomrule
\endlastfoot
$n\ge2$, $m=n-1$ & complex dimension of phase space; fibre rank of the toric models & \S\Mnum{sec:intro}\\
$z_j=x_j+iy_j$ & coordinates of $\C^n=\R^{2n}$ & \S\Mnum{sec:intro}\\
$\omega_0$, $X_H$ & $\sum dx_j\wedge dy_j$; Hamiltonian field, $\iota_{X_H}\omega_0=dH$ & \S\Mnum{sec:intro}\\
$J_0$, $\varphi^t_H$ & matrix of $\omega_0$ (multiplication by $-i$); flow of $X_H$ & \S\Mnum{sec:ham}, \S\ref{sec:pre}\\
$B(c)$, $\Bo(c)$, $B(s,c)$ & closed / open ball $\{\pi|z|^2\le c\}$ of capacity $c$; $s+B(c)$ & \S\Mnum{sec:intro}\\
$\EB(r)$, $\EB^\circ(r)$ & closed / open Euclidean ball of radius $r$ & \S\Mnum{sec:ham}\\
$k$, $c_i$, $s_i$, $C$ & number of balls, capacities, centres, target capacity & \S\Mnum{sec:intro}\\
$\rho_i=c_i/C$ & normalised capacities & \S\Mnum{sec:intro}\\
$a_i$; $\bar\rho_i$, $\hat\rho_i$ & auxiliary capacities; closed enlarged source capacity and open domain capacity of the layer construction, $\rho_i<\bar\rho_i<\hat\rho_i<a_i$; with $\delta_i=\frac12(a_i-\rho_i)$, $(\bar\rho_i,\hat\rho_i)=(a_i-\delta_i,a_i-\frac34\delta_i)$ for $n=2$ and $(a_i-\frac34\delta_i,a_i-\frac12\delta_i)$ for $n\ge3$ & \S\Mnum{sec:hdl}, Def.~\ref{def:str-enlarged}\\
$\eta_\star=1-\rho\sqrt k$ & volume defect of the equal-capacity branch in dimension four; $t_0=1-\eta_\star/2$ & \S\ref{sec:fcp-equal}, Rem.~\ref{rem:rtc-equal}\\
$\rho_i^+$, $\rho_i^{\rm dom}$ & enlarged closed source capacity and capacity of the open embedding domain, chosen separately in each branch & \S\Mnum{sec:ham}\\
$A$, $a$, $s=1-a$ & large given capacity, reserved capacity, shell parameter & \S\ref{sec:lrg}\\
$p=(p_j)$, $|p|$, $\theta_j$ & moments $\pi|w_j|^2$, their sum, angles $\arg w_j/2\pi$ & \S\Mnum{sec:intro}, \S\ref{sec:pre}\\
$\omega_V$, $X_{p_j}$ & fibre form $\sum dp_j\wedge d\theta_j$; rotation generator, $\iota_{X_{p_j}}\omega_V=-dp_j$ & \S\Mnum{sec:mech}, \S\ref{sec:pre}\\
$w=(w_j)$; $\varpi$ & fibre coordinates; moment of the base disc in $\C\times\C^m$ & \S\Mnum{sec:mech}\\
$\Delta$, $\langle b_\nu,p\rangle\le\ell_\nu$, $\tau_\Delta$ & downward moment polytope, its outer faces, core fraction & \S\S\Mnum{sec:mech}, \Mnum{sec:cmp}\\
$h$; $h_1,h_2$ & fibre height(s) of $\Delta$ & \S\Mnum{sec:mech}\\
$V_\Delta$, $V_\Delta(E)$ & toric domain $\{p(w)\in\Delta\}\subset\C^m$; its fibrewise version in $E$ & \S\Mnum{sec:mech}\\
$\Gamma$, $Z$ & base of a model and its marked set; the model section fixes the affine or cover presentation & \S\S\ref{sec:cub}--\ref{sec:lrg}\\
$\Sigma$ & base surface ($\Gamma$, or a compact bordered subsurface) & \S\Mnum{sec:mech}\\
$d^c$, $\omega_{\rm FS}$, $\sigma$ & $\frac i{4\pi}(\bar\partial-\partial)$; normalised affine Fubini--Study form; base area form & \S\ref{sec:pre}\\
$E$, $\alpha_j$, $\Omega_0$ & model space in its affine or cover presentation; angular forms; $\sigma+\sum d(p_j\alpha_j)$ & \S\S\ref{sec:cub}--\ref{sec:lrg}\\
$a_j$, $a_j(p)$, $a'_j$; $F_j$ & connection one-forms, toric mixed terms (not the capacities $a_i$); curvature (Chern) forms & \S\Mnum{sec:mech}, \S\ref{sec:pre}\\
$\kappa(p)$ & horizontal area form at moment $p$ & \S\Mnum{sec:mech}\\
$\Hb(p)=\int_\Sigma\kappa(p)$ & base area profile & \S\Mnum{sec:mech}\\
$\Pb(p)$ & unused area $\Hb(p)-\sum_i(a_i-|p|)_+$ & \S\Mnum{sec:mech}\\
$P$ & minorant of $\Pb$ in a comparison datum & \S\S\Mnum{sec:mech}, \Mnum{sec:cmp}\\
$K$, $F=K-P$, $\gamma$ & comparison function, comparison defect, gap & \S\S\Mnum{sec:mech}, \Mnum{sec:cmp}\\
$R_j$, $\Rcal$ & coordinatewise increasing rearrangement, $\Rcal=R_m\circ\dots\circ R_1$ & \S\Mnum{sec:cmp}\\
$q_t$, $q=q_1$, $Q_t$ & fibre isotopy, its endpoint, generator ($\iota_{X_{Q_t}}\omega_V=-dQ_t$) & \S\S\Mnum{sec:mech}, \Mnum{sec:cmp}\\
$\mathcal A\subset\mathcal A^+$, $\mathcal B$, $(s,v)$, $\beta$ & handle annulus, dual band, handle coordinates, closed one-form of period $1$ & \S\Mnum{sec:hdl}\\
$\nu$, $c_\nu$; $\Lambda$ & periodic unit-area reference form and its local density; finite affine reference form of the handle on a bordered base & \S\Mnum{sec:hdl}\\
$e_\Lambda\Lambda$, $H^*$, $J$ & positive background form, handle mass $\Hb-\epsilon_{\rm sm}$, collar coefficient $K-H^*$ & \S\Mnum{sec:hdl}\\
$h_i(p)$, $R^*$ & smooth caps, $0\le h_i\le(a_i-|p|)_+$, with $h_i\ge(a_i+\rho_i^\sharp)/2-|p|$; residual & \S\Mnum{sec:hdl}\\
$\mathcal D$, $\Phi$; $\Lambda_i$, $E_i$ & twist, Moser map; layers, layer maps ($\mathcal E_i$ in the main text) & \S\Mnum{sec:hdl}\\
$\mathcal T$, $\mathcal C_{\mathcal K}$ & transfer from a model to the target, compact set containing the image of $\mathcal K$ & \S\Mnum{sec:mech}, \S\S\ref{sec:cub}--\ref{sec:lrg}\\
$S$, $s_1$, $\sigma_j$; $\chi_{r,r'}$, $\chi^{\rm I}$ & flat step, $\max S'$, clocks; radial and interval cutoffs & \S\ref{sec:pre}\\
$\mathfrak v_t$ & Moser primitive, $\iota_{Y_t}\Omega_t=-\mathfrak v_t$ & \S\ref{sec:pre}\\
$\defdesc(h)$ & definite description & \S\Mnum{sec:intro}\\
$r_{1,i}<r_{2,i}<r_{3,i}<r_{4,i}$ & enlarged Euclidean source radii for the common isotopy and its cutoff & \S\Mnum{sec:ham}\\
$y_i=f_i(0)$, $Y_i=y_i/\tau$, $L$, $\delta_c$, $\tau$ & isotopy centres, far centres, Lipschitz bound, separation, common tail scale & \S\Mnum{sec:ham}\\
$F_{i,t}$ & common isotopy & \S\Mnum{sec:ham}\\
$\mathfrak a_{i,t}$, $\mathfrak h_{i,t}$ & source one-form, source primitive & \S\Mnum{sec:ham}\\
$H=H[n,k,C,c,s]$ & the global Hamiltonian & \S\Mnum{sec:ham}\\
$\beta_B$, $P$, $\Hfin$ & box cutoff, rational polynomial (this meaning of $P$ only in \S\Mnum{sec:eff}), $\Hfin=\beta_B\cdot P$ & \S\Mnum{sec:eff}\\
$b_R$, $L_X$, $\eta$ & radial budget, Lipschitz bound of $X_H$, field error (other uses of $\eta$, such as the connection form of the cubic model, are local) & \S\Mnum{sec:eff}\\
L1, L2, L3 & levels of explicitness & \S\Mnum{sec:intro}\\
\end{longtable}
}
